Cell morphology and gene expression capture complementary responses to perturbations, but differences in measurement modality and experimental context make mechanism-related cellular responses difficult to compare directly. We present BioMoAX, a multimodal retrieval model trained on Cell Painting images and L1000 expression profiles. Separate projection heads distinguish compound identity from mechanism-of-action (MOA) relationships, with five complementary contrastive objectives using paired measurements, shared MOA annotations and gene--drug associations. The same trained encoders and drug reference library support functional retrieval across measurement modalities, experimental batches, cellular backgrounds and perturbation types. HT29 images excluded from image training achieve Top-5 MOA hit rates of 21.78\% for image-to-image retrieval and 18.69\% for image-to-expression retrieval, compared with 8.45\% and 5.35\% for the strongest corresponding baselines, CellProfiler and ridge-mapped CLOOME. CRISPR image queries for held-out genes achieve a Top-10 hit rate of 19.77\%, compared with 13.79\% for CellProfiler. Two literature-supported co-retrieval cases connect functionally related genes through shared drugs, illustrating how phenotypic references complement incomplete annotations. BioMoAX enables images of drug-treated cells to query existing morphological and expression profiles and CRISPR images to retrieve functionally related drugs.
